
\documentclass[11pt]{article}
\usepackage[margin=1in]{geometry}
\usepackage{amsmath,amssymb,booktabs,array,longtable}
\usepackage{graphicx}
\usepackage{hyperref}
\usepackage{enumitem}
\usepackage{titlesec}

\begin{document}

\begin{center}
\Large\textbf{Shiv Nadar University Chennai}\\[2mm]
\end{center}

\renewcommand{\arraystretch}{1.4}

\begin{tabular}{|p{3.5cm}|p{8cm}|p{2cm}|}
\hline
Degree \& Branch & B.Tech Artificial Intelligence \& Data Science & Semester V\\ \hline
Subject Code \& Name & CS3807 -- Deep Learning Laboratory & AY: 2026--27\\ \hline
\end{tabular}

\vspace{3mm}

\begin{center}
\Large\textbf{Experiment 1}\\
\large\textbf{Implementation of a Single Layer Perceptron for Binary Classification}
\end{center}

%----------------------------------------------------------

\section*{1. Objective}

The objective of this experiment is to understand the concept of an artificial neuron and implement a Single Layer Perceptron from scratch for binary classification. Students will learn the perceptron learning algorithm, understand the importance of activation functions, visualize the learning process, and evaluate the classifier using a real-world dataset.

%----------------------------------------------------------

\section*{2. Tasks to be Performed}

\begin{enumerate}[leftmargin=*]
\item Download and understand the dataset.
\item Perform Exploratory Data Analysis (EDA).
\item Preprocess the dataset.
\item Implement a Single Layer Perceptron from scratch.
\item Train the model using the perceptron learning rule.
\item Evaluate the classifier using standard performance metrics.
\item Analyze the effect of different learning rates.
\item Compare the implementation with Scikit-learn's Perceptron (optional).
\end{enumerate}

%----------------------------------------------------------

\section*{3. Background Theory}

An Artificial Neuron is the fundamental building block of a neural network. It receives one or more input features, multiplies them by corresponding weights, adds a bias term, and applies an activation function to determine the output.

The Single Layer Perceptron is the earliest supervised learning algorithm capable of solving linearly separable binary classification problems.

\subsection*{Mathematical Formulation}

Weighted Sum

\[
z=\mathbf{w}^{T}\mathbf{x}+b
\]

where

\begin{itemize}
\item $\mathbf{x}$ : Input vector
\item $\mathbf{w}$ : Weight vector
\item $b$ : Bias
\item $z$ : Linear combination
\end{itemize}

Prediction

\[
\hat y=f(z)
\]

Weight Update Rule

\[
\mathbf{w}_{new}
=
\mathbf{w}_{old}
+
\eta
(y-\hat y)
\mathbf{x}
\]

Bias Update

\[
b_{new}
=
b_{old}
+
\eta
(y-\hat y)
\]

where $\eta$ denotes the learning rate.

%----------------------------------------------------------

\section*{4. Activation Functions}

An activation function determines the output of a neuron based on the weighted sum of its inputs. It introduces non-linearity, enabling neural networks to learn complex patterns. Although the perceptron uses only the Step Activation Function, modern deep learning models employ differentiable activation functions for efficient optimization through backpropagation.

\subsection*{Step Activation Function}

\[
f(z)=
\begin{cases}
1,&z\ge0\\
0,&z<0
\end{cases}
\]

The Step Activation Function produces a binary output and is therefore suitable only for linearly separable binary classification problems.

\begin{center}
\begin{tabular}{|p{2.5cm}|p{3cm}|p{6cm}|}
\hline
\textbf{Activation} & \textbf{Output Range} & \textbf{Typical Applications}\\
\hline
Step & $\{0,1\}$ & Classical Perceptron\\
\hline
Sigmoid & $(0,1)$ & Binary Classification Output Layer\\
\hline
Tanh & $(-1,1)$ & Hidden Layers (Traditional Networks)\\
\hline
ReLU & $[0,\infty)$ & Hidden Layers in CNNs and MLPs\\
\hline
Leaky ReLU & $(-\infty,\infty)$ & Avoiding Dying ReLU\\
\hline
Softmax & $(0,1)$ & Multi-class Classification Output Layer\\
\hline
\end{tabular}
\end{center}

\textbf{Note:} This experiment uses only the Step Activation Function. The remaining activation functions will be studied in subsequent experiments involving Multilayer Perceptrons and Deep Neural Networks.

%----------------------------------------------------------

\section*{5. Dataset Description}

\textbf{Dataset:} Banknote Authentication Dataset

\textbf{Source:} UCI Machine Learning Repository

\textbf{Download Link:}

\url{https://archive.ics.uci.edu/dataset/267/banknote+authentication}

\begin{center}
\begin{tabular}{|l|l|}
\hline
Instances & 1372\\
\hline
Features & 4 Numerical Features\\
\hline
Classes & 2\\
\hline
Missing Values & None\\
\hline
Task & Binary Classification\\
\hline
\end{tabular}
\end{center}

\textbf{Feature Description}

\begin{itemize}
\item Variance
\item Skewness
\item Curtosis
\item Entropy
\end{itemize}

Target Class

\begin{itemize}
\item 0 -- Authentic Banknote
\item 1 -- Forged Banknote
\end{itemize}

%----------------------------------------------------------

\section*{6. Experimental Procedure}

\textbf{Task 1: Dataset Exploration}

\begin{itemize}
\item Load the dataset.
\item Display the first five samples.
\item Determine dataset dimensions.
\item Identify missing values.
\item Display descriptive statistics.
\end{itemize}

\textbf{Expected Output}

Dataset summary and statistical information.

\vspace{2mm}

\textbf{Task 2: Exploratory Data Analysis}

Generate

\begin{itemize}
\item Histograms
\item Correlation Heatmap
\item Scatter Plot
\item Boxplots
\end{itemize}

\vspace{2mm}

\textbf{Task 3: Data Preprocessing}

\begin{itemize}
\item Normalize all numerical features.
\item Split the dataset into Training (80\%) and Testing (20\%).
\end{itemize}

\vspace{2mm}

\textbf{Task 4: Perceptron Implementation}

Implement from scratch

\begin{itemize}
\item Weight Initialization
\item Bias Initialization
\item Step Activation Function
\item Forward Propagation
\item Perceptron Learning Rule
\end{itemize}

\vspace{2mm}

\textbf{Task 5: Model Training}

Display

\begin{itemize}
\item Epoch Number
\item Misclassified Samples
\item Updated Weights
\item Updated Bias
\end{itemize}

\vspace{2mm}

\textbf{Task 6: Model Evaluation}

Compute

\begin{itemize}
\item Accuracy
\item Precision
\item Recall
\item F1-score
\item Confusion Matrix
\end{itemize}

%----------------------------------------------------------

\section*{7. Mandatory Plots}

\begin{longtable}{|p{3.4cm}|p{3.4cm}|p{6cm}|}
\hline
\textbf{Plot} & \textbf{Purpose} & \textbf{Expected Inference}\\
\hline
Feature Histograms & Distribution Analysis & Identify skewness and outliers.\\
\hline
Correlation Heatmap & Feature Relationship & Observe highly correlated features.\\
\hline
Scatter Plot & Class Distribution & Determine whether classes are approximately linearly separable.\\
\hline
Training Error vs Epoch & Learning Behaviour & Verify convergence of the perceptron.\\
\hline
Weight Evolution & Parameter Analysis & Observe stabilization of learned weights.\\
\hline
Bias Evolution & Parameter Analysis & Understand the role of the bias term.\\
\hline
Confusion Matrix & Performance Evaluation & Interpret TP, FP, TN and FN.\\
\hline
Learning Rate Comparison & Hyperparameter Study & Compare convergence behaviour for different learning rates.\\
\hline
Decision Boundary (Optional) & Visualization & Illustrate the separating hyperplane using two selected features.\\
\hline
\end{longtable}

Students should provide a brief interpretation (2--3 sentences) for every generated plot.

%----------------------------------------------------------

\section*{8. Performance Tables}

\textbf{Training Summary}

\begin{tabular}{|p{6cm}|p{7cm}|}
\hline
Dataset Size &\\
\hline
Train/Test Split &\\
\hline
Learning Rate &\\
\hline
Epochs &\\
\hline
Final Weights &\\
\hline
Final Bias &\\
\hline
Accuracy &\\
\hline
Precision &\\
\hline
Recall &\\
\hline
F1-score &\\
\hline
\end{tabular}

\vspace{3mm}

\textbf{Epoch-wise Learning}

\begin{tabular}{|c|c|c|c|c|}
\hline
Epoch & Errors & Weight 1 & Weight 2 & Bias\\
\hline
&&&&\\
&&&&\\
&&&&\\
\hline
\end{tabular}

%----------------------------------------------------------

\section*{9. Additional Tasks}

\begin{enumerate}[leftmargin=*]
\item Compare the Step and Sigmoid activation functions. Discuss why the Step Activation Function is not used in modern deep learning models.
\item Compare the implementation with Scikit-learn's Perceptron.
\item Repeat the experiment using learning rates 0.001, 0.01 and 0.1.
\item Explain why a Single Layer Perceptron cannot solve the XOR problem.
\item Study the effect of feature normalization on model convergence.
\end{enumerate}

%----------------------------------------------------------

\section*{10. Report Contents}

The report should include:

\begin{itemize}
\item Objective
\item Background Theory
\item Activation Functions
\item Dataset Description
\item Experimental Procedure
\item Source Code
\item Experimental Results
\item Generated Plots with Interpretation
\item Performance Tables
\item Discussion
\item Conclusion
\item References
\end{itemize}

%----------------------------------------------------------

\section*{11. References}

\begin{enumerate}
\item F. Rosenblatt, ``The Perceptron,'' \emph{Psychological Review}, 1958.
\item I. Goodfellow, Y. Bengio and A. Courville, \emph{Deep Learning}, MIT Press, 2016.
\item C. M. Bishop, \emph{Pattern Recognition and Machine Learning}, Springer, 2006.
\item S. Haykin, \emph{Neural Networks and Learning Machines}, Pearson, 2009.
\item UCI Machine Learning Repository -- Banknote Authentication Dataset.
\item Scikit-learn Documentation: Perceptron.
\end{enumerate}

\end{document}

